\chapter{Sample Playback}
\label{ch:playback}

Each of the six voices is an independent player. It holds a position in the sample and a step, reads one byte of the
sample for every output sample, and adds the step to the position. This chapter describes how a voice starts a note,
how it plays and loops a sample, and how its pitch is set, adjusted and glided.

\section{Voice State}

A voice keeps these quantities, none of which the CPU can read directly:

\begin{itemize}
	\item The \textbf{instrument record} of the instrument it last loaded (Section~\ref{sec:instrument-table}).
	\item The \textbf{position} \texttt{SPA}, a 16.16 number. Its integer part is the byte offset into the sample.
	\item The \textbf{step} \texttt{SPI}, a 16.16 number: the sample bytes the position advances for each output sample.
	\item The \textbf{target step}, which the step moves towards when the voice glides (Section~\ref{sec:portamento}).
	\item The \textbf{envelope level} and phase (Chapter~\ref{ch:envelope}).
\end{itemize}

\section{Starting a Note}
\label{sec:note-start}

A note starts when the sequencer executes \texttt{NOTE} or the CPU writes \texttt{COMMAND} with bit 15 set. Both give a
note number from 0 to 119, and both use the instrument number in the voice's \texttt{INSTR} register. Starting a note
takes these steps, in order:

\begin{enumerate}
	\item \textbf{Read the instrument.} If the instrument number differs from the one the voice last loaded, or the voice has
	      loaded none, the device reads the instrument record: words 0 to 6 of Table~\ref{tab:instrument}. A note on the
	      instrument the voice has already loaded does not read the record again.
	\item \textbf{Read the pitch.} The device reads the two words of the pitch table entry for the semitone,
	      \texttt{SPIT\_ADDRESS} $+ 4 (n \bmod 12)$, and finds the step as in Section~\ref{sec:note-numbers}. The signed
	      \texttt{FINE} value is added to it (Section~\ref{sec:fine-pitch}).
	\item \textbf{Apply.} The voice either \emph{retriggers} or \emph{glides}. It glides when \texttt{PORTA\_EN} is set, the
	      instrument is the one the voice is already playing, and the voice is still sounding, meaning its envelope is not
	      inactive. A glide changes only the target step. In every other case the voice retriggers: the step and the target
	      step are both set to the new step, the position is set to 0, the envelope restarts, and the first envelope step is taken
	      at once (Chapter~\ref{ch:envelope}).
\end{enumerate}

A change of instrument always retriggers, whether or not portamento is enabled. Reading the instrument and the pitch
entry takes up to nine reads of audio memory, which are made in the sound processor's memory slots after the sample
reads (Chapter~\ref{ch:timing}); a note therefore takes effect a little after it is requested, and the CPU is not told
when. A note that is requested while an earlier one is still being applied is applied after it.

The device reads the instrument record when a voice's instrument changes and not on every note. A program that changes
the contents of an instrument record that a voice has loaded must start a note on a different instrument on that voice
before the change is seen.

\section{Playing a Sample}
\label{sec:playing}

For every output sample, each voice does the following with its position and step:

\begin{enumerate}
	\item It reads the sample byte at the integer part of the position, relative to \texttt{SAMPLE\_START}, as a signed 8-bit
	      value. The first byte read after a retrigger is byte 0.
	\item It adds the step to the position.
	\item It applies the loop rules of Section~\ref{sec:looping} to the new position.
\end{enumerate}

The voice does not interpolate: the byte at the integer position is used as it is, and the fraction of the position
only keeps the pitch accurate over time. A step of 1.0 plays the sample at the output rate. A step of 2.0 plays it an
octave higher and skips every second byte; a step of 0.5 plays it an octave lower and repeats every byte.

The nearest-sample lookup adds error to the sound that depends on the length of the waveform and not on the pitch. A
waveform of 64 bytes played this way carries distortion about 24 dB below the signal; a waveform of 256 bytes, about 36
dB. The error falls by about 6 dB for each doubling of the length, so a long waveform sounds cleaner than a short one.
Chapter~\ref{ch:programming} gives guidance.

\subsection{Looping}
\label{sec:looping}

When the integer part of the position reaches \texttt{SAMPLE\_LENGTH} or more:

\begin{itemize}
	\item With \texttt{LOOP\_ENABLE} set, the device subtracts \texttt{SAMPLE\_LENGTH} $-$ \texttt{LOOP\_START} from the integer
	      part, once, and keeps the fraction. The overshoot is preserved, so the pitch does not drift. The sample then plays from
	      \texttt{LOOP\_START} for ever.
	\item With \texttt{LOOP\_ENABLE} clear, the sample has ended. The voice's output is zero from the next output sample, until a
	      new note retriggers it. The envelope keeps running, so a voice whose sample has ended is silent, not stopped.
\end{itemize}

The device subtracts the loop length only once for each output sample. A step that is larger than the loop length
(\texttt{SAMPLE\_LENGTH} $-$ \texttt{LOOP\_START}) is architecturally undefined. A waveform that holds one cycle and is
played at any pitch below half the output rate never meets this limit, because its step is less than half of its
length.

\section{Pitch Adjustment}
\label{sec:fine-pitch}

The \texttt{FINE} register holds a signed 13-bit value, $-4096$ to $+4095$, that is added to the step as raw 16.16
units when a note is started. It is added after the octave shift, so a unit is the same size for every note of the
instrument and not the same fraction of its pitch. Table~\ref{tab:fine-units} shows what that means for the 64-byte
waveform of Table~\ref{tab:pitch-table}: a unit moves the pitch of a low note by a fraction of a cent that is large and
of a high note by a fraction that is very small.

\begin{table}[H]
	\centering
	\begin{tabular}{rrrr}
		\toprule
		\textbf{Note} & \textbf{Step (16.16)} & \textbf{One unit (cents)} & \textbf{$+$4095 units (cents)} \\
		\midrule
		0             & 3,609                 & 0.480                     & $+1313$                        \\
		12            & 7,218                 & 0.240                     & $+778$                         \\
		24            & 14,436                & 0.120                     & $+432$                         \\
		48            & 57,744                & 0.030                     & $+119$                         \\
		96            & 923,904               & 0.0019                    & $+7.7$                         \\
		119           & 3,488,256             & 0.0005                    & $+2.0$                         \\
		\bottomrule
	\end{tabular}
	\caption{The Size of a Fine Pitch Unit (64-Byte Waveform)}
	\label{tab:fine-units}
\end{table}

A program that detunes by a fixed musical amount across octaves scales the value it writes by the octave, or writes
\texttt{FINE} before each note.

\section{Portamento}
\label{sec:portamento}

When a voice glides, its step moves from the current step to the target step, a little at a time, and the pitch slides
with it. \texttt{PORTA\_RATE} is the size of each move, in raw units of the 16.16 step. The voice makes one move for
every envelope update, which is once in every 64 output samples, about every 2.0 ms (Chapter~\ref{ch:envelope}). The
step moves up or down towards the target, and a move that would pass the target stops at it. A rate of zero holds the
step where it is until another note changes it.

Portamento runs on the voice's own envelope clock, so it needs neither the sequencer nor a particular tick rate: it
runs when \texttt{SEQ\_PERIOD} is 0, and for a voice the CPU controls directly. It also runs for a voice whose
\texttt{PORTA\_EN} is clear, if the voice has a target that differs from its step; setting \texttt{PORTA\_EN} only
decides whether a new note glides or retriggers.

The glide time is the distance between the two steps divided by the rate, times 2.0 ms. The step is linear in
frequency, so the slide is linear in Hz and not in semitones. The slowest rate, 1, moves the step by one raw unit in
2.0 ms, so the longest glide between two notes is proportional to the distance between them: an octave up from A0 in
Table~\ref{tab:pitch-table} is 3,609 units and takes at most 7.2 s. A short glide at a low pitch is limited the other
way: a semitone up from A0 is 215 units and takes at most 0.43 s.

\section{Stopping a Note}

A stop, whether the sequencer's \texttt{STOP} or a \texttt{COMMAND} with bit 15 clear, starts the voice's envelope
release. The voice carries on playing its sample under the falling envelope, so the sound dies away and is not cut
(Chapter~\ref{ch:envelope}). A stop on a voice whose envelope is inactive has no effect.
